




\subsection{5.1, Eigenvalues and eigenvectors}

\begin{problems}
%
\problem % Jankowski Spring 2018 exam 3
True or false: If $A$ is an $n \times n$ matrix and $\det(A)=2$, then $2$ is an
eigenvalue of $A$.

\begin{solution}
  False.  For example, $A=\mat{4 0; 0 1/2}$ has $\det(A)=2$ but its eigenvalues
  are $4$ and $\dfrac{1}{2}$.
\end{solution}

\problem % Created for ILA
True or false: Suppose $A$ is a $n \times n$ matrix and $u$ and $v$ are nonzero
vectors in ${\bf R}^{n}$ satisfying $Au = 2u$ and $Av = 3v$.  Then $\{u, v\}$ must
be linearly independent.

\begin{solution}
  True: we have implicitly been told that $2$ and $3$ are eigenvalues of $A$,
  where $u$ is a $2$-eigenvector and $v$ is a $3$-eigenvector.  Since eigenvectors
  corresponding to different eigenvalues are automatically linearly independent,
  the set $\{u,v\}$ must be linearly independent.
\end{solution}

\problem % Jankowski Spring 2018 exam 3
True or false: The matrices $A=\mat{2 1; 0 -1}$ and $B=\mat{2 20; 0 -1}$ have
the same eigenvalues.
    
\begin{solution}
  True.  Each matrix has eigenvalues $\lambda=2$ and $\lambda = -1$.
\end{solution}


\problem % Jankowski Spring 2018 exam 3
True or false: If $A$ is an $n \times n$ matrix and $v$ and $w$ are eigenvectors
of $A$, then $v+w$ is also an eigenvector of $A$.
    
\begin{solution}
  False.  For example, if $A= \mat{1 0; 0 2}$ then $\vec{1 0}$ and $\vec{0 1}$
  are eigenvectors, but $A \vec{1 1} = \vec{1 2}$ so $\vec{1 1}$ is not an
  eigenvector.
\end{solution}

\problem % Dan's notes
True or false: The zero vector is an eigenvector for every square matrix.
    
\begin{solution}
  False.  By definition of eigenvector, the zero vector is never an eigenvector
  for any matrix.
\end{solution}

\problem % Created for ILA
Which of the following vectors (if any) are eigenvectors for
$A = \mat{35 -32; 40 -37}$?
% eigenvalues $3$ with e.v. (1,1)... and eigenvalue -5 with e.v. (4,5)
\begin{subproblems}
  \subproblem  $v_1 = \vec{2 1}$  %
  \subproblem  $v_2 = \vec{-1 -1}$%
  \subproblem  $v_3 = \vec{4 5}$ %
\end{subproblems}

\begin{solution}
We check each vector one-by-one.
\[Av_1 = \mat{35 -32; 40 -37} \vec{2 1} = \vec{38 43}, %
  \quad \textup{which is not a scalar multiple of } v_1.\]
\[Av_2 = \mat{35 -32; 40 -37} \vec{-1 -1} = \vec{-3 -3} = 3v_2.\]
\[Av_3 = \mat{35 -32; 40 -37} \vec{4 5} = \vec{-20 -25} = -5v_3.\] From the
computations, we conclude that $v_2$ and $v_3$ are eigenvectors of $A$, but $v_1$
is not an eigenvector of $A$ because $Av_1$ is not a scalar multiple of $v_1$.
\end{solution}


\problem %Dan's notes 
Find the $2$-eigenspace for the matrix
\[\mat{0 13 12; 1/4 0 0; 0 1/2 0}.\]
\begin{solution}
We need to find $\Nul(A-2I)$.  For this, we note
\[A - 2I = \mat{0 13 12; 1/4 0 0; 0 1/2 0} - \mat{2 0 0; 0 2 0; 0 0 2}
= \mat{-2 13 12; 1/4 -2 0; 0 1/2 -2}.\]
Now:
\[\amat{A-2I 0} =
  \amat{-2 13 12 0; 1/4 -2 0 0; 0 1/2 -2 0}
  \xrightarrow[]{RREF}
  \amat{1 0 -32 0; 0 1 -4 0; 0 0 0 0}.\]
This gives us $x_1 - 32x_3 = 0$ and $x_2 - 4 x_3 = 0$ with $x_3$ free.  The parametric
vector form is 
\[\vec{x_1 x_2 x_3} = \vec{32x_3 4x_3 x_3} = x_3 \vec{32 4 1}.\]
Therefore, the $2$-eigenspace of $A$ is
\[\Nul(A-2I) = \Span \left\{\vec{32 4 1}\right\}.\]
\end{solution}

\problem %Dan's notes
How many different eigenvalues can an $n \times n$ matrix have?

\begin{solution}
  The maximum is $n$.  Since eigenvectors for different eigenvalues are
  automatically linearly independent, and since no set of vectors in
  ${\bf R}^{n}$ can be linearly independent if it has more than $n$ vectors in
  it, we cannot have more than $n$ different eigenvalues.  It is possible to
  have exactly $n$ different eigenvalues, for example the $3 \times 3$ matrix
  below has $3$ different eigenvalues.
  \[\mat{1 0 0; 0 2 0; 0 0 3}.\]
\end{solution}

\problem %Dan's notes
Consider the matrix $A$ for reflection of ${\bf R}^2$ about the line $y = 7x$.
What are its eigenvalues?

\begin{solution}
  The eigenvalues are $\lambda = 1$ and $\lambda = -1$.  The matrix does not
  change any vector along the line $y = 7x$.  In other words, $Av = v$ for every
  vector $v$ along the line $y = 7x$. Therefore, $\lambda = 1$ is an eigenvalue.
  \newline \newline %
  The matrix flips every vector on the line through the origin that is
  perpendicular to the line $y = 7x$.  In other words, $Av = -v$ for every
  vector $v$ along the line $y = -x/7$.  Therefore, $\lambda = -1$ is an
  eigenvalue. \newline \newline %
  Since a $2 \times 2$ matrix can have at most $2$ eigenvalues, we conclude that
  $1$ and $-1$ are the only eigenvalues of $A$. \newline \newline %
  Note that the eigenvalues and reasoning would be identical if we changed the
  line $y = 7x$ to a different line $L$ through the origin.  The eigenvalues
  would still be $1$ and $-1$, the $1$-eigenspace would be $L$, and the
  $(-1)$-eigenspace would be the line through the origin perpendicular to $L$.
  
\end{solution}


\problem
Without doing any major computations, write one real eigenvalue for
each of the following matrices. 
\begin{subproblems}
  \subproblem $A = \mat{6 0 0; -4 6 0; 6 6 6}$ \subproblem The
  standard matrix $A$ for the linear transformation $T$ that satisfies
  \[T(e_1) = e_3, \quad T(e_2) = 2e_2 - e_3, \quad T(e_3) = -e_2
    +2e_3.\]
  %
  \subproblem
  $A = \mat{1 2 3 4 5 6; 1 2 3 4 5 6; 0 0 0 1 0 0; 0 0 0 0 1 0; 0 0 0
    0 0 1}.$
\end{subproblems}

\begin{solution}
  \begin{subproblems}
    \subproblem Since $A$ is lower-triangular, its eigenvalues are its
    diagonal entries.  Each of these is $6$, so the only eigenvalue of
    $A$ is $\lambda = 6$.
     %
    \subproblem We see $e_2+e_3$ is an eigenvector corresponding to
    eigenvalue $\lambda = 1$ since
    \[T(e_2+e_3) = 2e_2-e_3 + (-e_2 + 2e_3) = e_2 + e_3.\]
    %
    \subproblem The first two rows of the square matrix $A$ are the
    same, so without doing any work we know that $A$ is not invertible,
    therefore $\lambda = 0$ is an eigenvalue of $A$ since the equation
    $Ax = 0$ has a non-trivial solution.

    Alternatively, we could have noted from the form of $A$ that
    $A(e_1-e_2) = 0$, which shows that $\lambda = 0$ is an eigenvalue
    and $e_1 - e_2$ is a corresponding eigenvector.
   \end{subproblems}
 \end{solution}

 
\problem %Jankowski Spring 2018 exam 2
    Let $T: {\bf R}^2 \to {\bf R}^2$ be the linear transformation
    which reflects across the line $L$ drawn below, and let $A$ be the
    standard matrix for $T$.  Answer both questions below without
    doing any computations.  \newline
    \begin{center}
      \includegraphics[scale=0.5]{q6_18s-line.pdf}
    \end{center}

	\begin{subproblems} % Jankowski quiz from forgot when
      \subproblem Find all eigenvalues of $A$.
      %
      \subproblem For each eigenvalue of $A$, draw one eigenvector on
      the graph above.  Your eigenvector does not need to be perfect,
      but it should be reasonably accurate. \newline
    \end{subproblems}


    \begin{solution}
      \begin{subproblems}
        \subproblem The equation of the line was not given, and it is
        irrelevant.  The eigenvalues of $A$ are $\lambda_1 = 1$ and
        $\lambda_2 = -1$, since $A$ fixes every vector along the line
        $L$ while $A$ flips every vector perpendicular to $L$.

        In the graph, $v_1$ corresponds to $\lambda_1 = 1$, while
        $v_2$ corresponds to $\lambda_2= -1$.

        \subproblem We graph eigenvectors below.
        \begin{center}
          \includegraphics[scale=0.4]{q6_18s-answer}
        \end{center}

        Many answers are possible: $v_1$ can be any nonzero vector on
        $L$ (going up-to-right or down-to-left), while $v_2$ can be
        any nonzero vector perpendicular to $L$ (going up-to-left or
        down-to-right).  \newline \newline Note: Algebraically, the
        problem would have been much more painful! \newline The line
        is actually the line $y = \dfrac{3}{7} x$, the matrix is
        $A= \dfrac{1}{29} \mat{20 21; 21 -20}$, and the eigenvectors
        drawn are $v_1 = \vec{7/3 1}$ and $v_2 = \vec{-1 7/3}$.
      \end{subproblems}
    \end{solution}

\problem % Created for ILA
Suppose $A$ is a $2 \times 2$ matrix, and that $u$ and $v$ are vectors
in the $(-2)$-eigenspace of $A$ and the $5$-eigenspace of $A$, respectively.
Find $A(8u - 4v)$ in terms of $u$ and $v$.

\begin{solution}
We are given that $Au = -2u$ and $Av = 5v$, so we use linearity:
\[A(8u-4v) = A(8u) - A(4v) = 8Au - 4Av = -16u -20v.\]
\end{solution}

\problem %Created for ILA
There is one $2 \times 2$ matrix $A$ with the property that
\[\Nul(A-2I) = \Span \left\{\vec{7 1}\right\},
  \quad \textup{ and } \quad
  \Nul(A+I) = \Span \left\{\vec{-1 3} \right\}.\]
Find $A$.

\begin{solution}
  We are told the $2$-eigenspace of $A$ is spanned by $\vec{7 1}$ and that the
  $(-1)$-eigenspace of $A$ is spanned by $\vec{1 -3}$, so
  \[Au = 2u = \vec{14 2} \quad \textup{ and } \quad
  Av = -v = \vec{1 -3}.\]
  Recall from section 3.3 that if $M$ is any $2 \times 2$ matrix, then its first
  column is $Me_1$ and its second column is $Me_2$ (this can also
  be seen directly from the definition of ``matrix times vector'' multiplication), so
  we need to find $A\vec{1 0}$ and $A \vec{0 1}$.

  For this, we write the standard basis vectors as linear combinations of the
  eigenvectors $u = \vec{7 1}$ and $v = \vec{1 -3}$. %
  \[\textup{For } \vec{1 0}: \amat{7 1 1; 1 -3 0}
    \xrightarrow[]{RREF} \amat{1 0 3/22; 0 1 1/22},
    \quad \textup{ so } \quad \dfrac{3}{22}u + \dfrac{1}{22} v
    = \vec{1 0},\]
  \begin{align*}
    A \vec{1 0}
    & = A \left(\dfrac{3}{22}u + \dfrac{1}{22} v\right) 
      = \dfrac{3}{22} Au + \dfrac{1}{22} Av 
      = \dfrac{3}{22} \vec{14 2} + \dfrac{1}{22} \vec{-1 3} \\
    &  = \vec{41/22 9/22}.
  \end{align*}

  \[\textup{For } \vec{0 1}: \amat{7 1 0; 1 -3 1}
    \xrightarrow[]{RREF} \amat{1 0 1/22; 0 1 -7/22},
    \quad \textup{ so } \quad \dfrac{1}{22}u - \dfrac{7}{22} v
    = \vec{0 1},\]
  \begin{align*}
    A \vec{0 1}
    & = A \left(\dfrac{1}{22}u - \dfrac{7}{22} v\right) 
      = \dfrac{1}{22} Au - \dfrac{7}{22} Av 
      = \dfrac{1}{22} \vec{14 2} - \dfrac{7}{22} \vec{-1 3} \\
    &  = \vec{21/22 -19/22}.
  \end{align*}

Therefore,
\[A = \mat{A\vec{1 0} A\vec{0 1}} = \mat{41/22 21/22; 9/22 -19/22}
  = \dfrac{1}{22} \mat{41 21; 9 -19}.\]

This problem's matrix $A$ is an example of a diagonalizable matrix.  We will
find a different way to compute $A$ in section 5.4.
\end{solution}




\end{problems}\setcounter{problem}{0}

\newpage

\subsection{5.2, The characteristic polynomial}
\begin{problems}
%
\problem % Dan's notes
True or false: Every $n \times n$ matrix has a real eigenvalue.

\begin{solution}
  False.  For example, the $2 \times 2$ matrix for rotation by $\pi/2$ radians
  counterclockwise has no real eigenvalues.  The matrix is $\mat{0 -1; 1 0}$,
  and we can compute that its characteristic equation is
  $0 = \det(A-\lambda I) = \lambda^2+1$, so its eigenvalues are $\pm i$.  We
  discuss complex eigenvalues in section 5.5.
\end{solution}

\problem % Dan's notes
True or false: Every $n \times n$ matrix has $n$ distinct eigenvalues.

\begin{solution}
False.  The $2 \times 2$ identity matrix $I = \mat{1 0; 0 1}$ has only one
eigenvalue, namely $\lambda = 1$.
\end{solution}

\problem % Dan's notes
True or false: If $\lambda$ is an eigenvalue of a matrix $A$, then the nullity
of $A- \lambda I$ is the dimension of the $\lambda$-eigenspace.

\begin{solution}
True.  By definition, the $\lambda$-eigenspace is the null space of 
$A-\lambda I$, so the dimension of the $\lambda$-eigenspace
is precisely the dimension of the null space of $A-\lambda I$.
\end{solution}


\problem %Dan's notes
What are the eigenvalues for the standard matrix for a reflection?

\begin{solution}
  They are always $1$ and $-1$.
\end{solution}

\problem %Dan's notes
What are the eigenvalues and eigenvectors for the $n \times n$ zero matrix?

\begin{solution}
  Let $A$ be the $n \times n$ zero matrix.  Then $A-\lambda I$ is the
  $n \times n$ diagonal matrix whose diagonal entries are all $-\lambda$, so its characteristic
  polynomial is $(-\lambda)^n$.  Therefore, the only eigenvalue is $\lambda = 0$.  The
  $0$-eigenspace is all of ${\bf R}^{n}$, since $Ax = 0$ for every $x$ in ${\bf R}^{n}$.
\end{solution}

\ifanswers \newpage \fi 

\problem %Dan's notes
Find the eigenvalues of 
$\mat{1 2 1; 0 -5 0; 1 8 0}.$
% Solution needed

\begin{solution}
We solve for $0 = \det(A-\lambda I)$ using the cofactor expansion
along the second row.
\begin{align*}
0 & = \det \mat{1-\lambda \, 2 1; 0 -5-\lambda \, 0; 1 8 -\lambda}
= (-1)^{4}(-5-\lambda) \det\mat{1-\lambda \, 1; 1 -\lambda}
\\ &
= (-5-\lambda)\Big(\lambda^2-\lambda-1\Big).
\end{align*}
Thus, $\lambda = -5$ is an eigenvalue, and so are the solutions to $\lambda^2-\lambda -1 = 0$.  For this, we use the quadratic formula to get
\[\lambda = \frac{-(-1) \pm \sqrt{(-1)^2-4(-1)(1)}}{2}
= \frac{1 \pm \sqrt{5}}{2}.\]
The eiegenvalues are $\lambda = -5$,
$\lambda = \dfrac{1 + \sqrt{5}}{2}$, and 
$\lambda = \dfrac{1 - \sqrt{5}}{2}$.
\end{solution}

\problem
Find the eigenvalues of $A = \mat{9 9; 9 9}$.  For each eigenvalue, find
a basis for the corresponding eigenspace.

\begin{solution}
  The characteristic polynomial can be computed by doing the full computation
  of $\det(A- \lambda I)$ or using the shortcut $2 \times 2$ formula.  It is
  \[\det(A-\lambda I) = \lambda^2 - \textup{Tr}(A) \lambda + \det(A)
    = \lambda^2 - 18 \lambda + 0 = \lambda(\lambda-18).\] Therefore,
  the eigenvalues are $\lambda = 0$ and $\lambda = 18$.  We row-reduce
  \[\amat{A-0I \ 0}= \amat{9 9 0; 9 9 0} \to \amat{1 1 0; 0 0 0},\]
  so $\left\{\vec{-1 1}\right\}$ is a basis for the $0$-eigenspace.
  Also,
  \[\amat{A-18I \ 0}= \amat{-9 9 0; 9 -9 0} \to \amat{1 -1 0; 0 0 0},\]
  so $\left\{\vec{1 1}\right\}$ is a basis for the $18$-eigenspace.
  
\end{solution}

\ifanswers \newpage \fi 

\problem %Dan's notes
Consider the matrix $A$ for rotation by $\pi/2$ radians about the $z$-axis in
${\bf R}^3$. What are its real eigenvalues?

\begin{solution}
  We can think of $A$ as the matrix that rotates the ``xy-piece'' of the vector
  by $\pi/2$ radians, while keeping the $z$-coordinate fixed.  This means
  \[A = \mat{0 -1 0; 1 0 0; 0 0 1}.\]
  Its characteristic polynomial is the following, which we compute using the cofactor
  expansion along the third row.
  \begin{align*}
    \det(A - \lambda I)
    & = \det\mat{-\lambda \, -1 0; 1 -\lambda \, 0; 0 0
      1-\lambda} = (-1)^{3+3}(1-\lambda)\Big[(-\lambda)(-\lambda)-1(-1)\Big]
    \\ & = (1-\lambda)(\lambda^2 + 1).
  \end{align*}
  The characteristic polynomial factors as $(1-\lambda)(\lambda+i)(\lambda-i)$,
  so the only real eigenvalue is $\lambda = 1$.  Complex eigenvalues will be
  covered in a later section of the textbook.
\end{solution}


\problem In each case, find the eigenvalues of the matrix $A$ and a basis
for each eigenspace.
\begin{subproblems}
  \subproblem $A$ the $2 \times 2$ matrix that reflects vectors across
  the line $y = x$.
  %
  \subproblem $A$ the $3 \times 3$ matrix that rotates vectors by
  $180^\circ$ around the $z$-axis, where we label points in
  ${\bf R}^3$ as $(x, y, z)$.
  %
\end{subproblems}

\begin{solution}
  \begin{subproblems}
    \subproblem The eigenvalues of $A$ are $\lambda = 1$ and
    $\lambda = -1$, since $A$ fixes vectors along the line $y=x$ and
    flips vectors perpendicular to the line $y=x$.
    %
    \subproblem The eigenvalues of $A$ are $\lambda = 1$ and
    $\lambda = -1$, since $A$ fixes vectors along the $z$-axis while
    $Ae_1 = -e_1$ and $Ae_2 = -e_2$.
  \end{subproblems}
\end{solution}

\problem Let $A$ be the $2 \times 2$ matrix satisfying
\[A\vec{1 3} = \vec{-2 1}, \qquad A \vec{0 1} = \vec{1 2}.\]
What are the eigenvalues of $A$?

\begin{solution}
  We find $A$ first.  Since $\vec{1 0} = \vec{1 3} - \vec{0 3}$,
  linearity gives us
  \[A\vec{1 0} = A \left(\vec{1 3} - \vec{0 3}\right) = A \vec{1 3} -
    3 A \vec{0 1} = \vec{-2 1} - 3 \vec{1 2} = \vec{-5 -5}.\]
  Therefore,
  \[A = \mat{A\vec{1 0} A \vec{0 1}} = \mat{-5 1; -5 2}.\]
  For the eigenvalues of $A$, we solve:
  \begin{align*}
    0 & = \det(A-\lambda I) = \lambda^2- \Tr(A)\lambda + \det(A)
    \\ & = \lambda^2 + 3 \lambda -5, 
    \\ & \lambda = \dfrac{-3 \pm \sqrt{9-4(-5)(1)}}{2}
         = \dfrac{-3 \pm \sqrt{29}}{2}.
  \end{align*}
\end{solution}

\problem Find the eigenvalues of the matrix $A$ that first reflects
vectors in ${\bf R}^2$ across the line $y = -x$, then rotates
$\pi/2$ radians counterclockwise, then scales the first coordinate
by a factor of $4$ while keeping the second coordinate fixed.

\begin{solution} One way to do this problem is to compute $A$.
  \begin{align*}
    A & = \mat{4 0; 0 1} \mat{0 -1; 1 0} \mat{0 -1; -1 0}
        = \mat{0 -4; 1 0} \mat{0 -1; -1 0} 
        = \mat{4 0; 0 -1}.
  \end{align*}
  Therefore, the eigenvalues of $A$ are $\lambda = 4$ and
  $\lambda = -1$.  \newline \newline Alternatively, we could have
  computed $Ae_1$ and $Ae_2$ geometrically and found that
  $Ae_1 = 4 e_1$ and $Ae_2 = -e_2$, so $e_1$ and $e_2$ are
  eigenvectors corresponding to the eigenvalues $4$ and $-1$
  respectively.  Since a $2 \times 2$ matrix has at most $2$ distinct
  eigenvalues, these are the only eigenvalues of $A$.

  
\end{solution}


\problem Suppose $A$ and $B$ are $n \times n$ matrices with the same
characteristic polynomial.  Do $A$ and $B$ have the same eigenvalues?
Answer sometimes, always, or never, and justify your answer.

\begin{solution}
  Always.  The eigenvalues of a square matrix are the roots of its
  characteristic polynomial, so if $A$ and $B$ have the same characteristic
  polynomial then their eigenvalues will be the same.
\end{solution}

\end{problems}\setcounter{problem}{0}

\newpage

\subsection{5.4, Diagonalization}

\begin{problems}
%

\problem % Dan's notes
True or false: If $A$ is a $3 \times 3$ matrix with eigenvalues $0$, $1$, and $2$, then
$A$ is diagonalizable.

\begin{solution}
  True: any $3 \times 3$ matrix with $3$ real eigenvalues is automatically
  diagonalizable.  One way to see this is that each of the $3$ eigenvalues must
  have a nonzero eigenvector in their eigenspace, and eigenvectors corresponding
  to different eigenvalues are linearly independent, therefore we obtain $3$
  linearly independent eigenvectors.  In other words, we obtain a basis for
  ${\bf R}^{3}$ consisting of eigenvectors of $A$, which is equivalent to the
  property that $A$ is diagonalizable.
\end{solution}


\problem % Dan's notes
True or false: If $\lambda$ is an eigenvalue of a matrix $A$, it is possible
that the $\lambda$-eigenspace is $0$-dimensional.

\begin{solution}
  False.  It is a fundamental fact that the geometric multiplicity of any
  eigenvalue $\lambda$ is at least $1$.  By definition of eigenvalue, the
  equation $Ax = \lambda x$ must have a non-trivial solution, therefore
  $(A-\lambda I)x = 0$ has a non-trivial solution, which means that
  $\Nul(A-\lambda I)$ is at least $1$-dimensional.
\end{solution}

\problem % Dan's notes
True or false: All diagonalizable matrices are invertible.

\begin{solution}
  False: for example, zero matrix is diagonalizable (in fact, diagonal!) but is
  not invertible.
\end{solution}

\problem % Dan's notes
True or false: All invertible matrices are diagonalizable.   

\begin{solution}
  False: for example, take $A = \mat{1 1; 0 1}$. Note $\det(A) = 1$, so $A$ is
  invertible.  However, the only eigenvalue of $A$ is $\lambda = 1$, and we can
  compute that the $1$-eigenspace is only $1$-dimensional, so $A$ is not
  diagonalizable.
\end{solution}

\problem % Dan's notes
True or false: All diagonal $n \times n$ matrices are diagonalizable.

\begin{solution}
  True: if $A$ is diagonal, then $A = I A I^{-1}$, so $A$ is diagonalizable.  In
  other words, $A = CDC^{-1}$ where $C$ is the $n \times n$ identity matrix and
  $D$ is $A$ itself!
\end{solution}

\problem % Dan's notes
True or false: All upper-triangular $n \times n$ matrices are diagonalizable.

\begin{solution}
False: for example,
\[A = \mat{2 1; 0 2}.\] The only eigenvalue of $A$ is $\lambda = 2$, but we can
compute that the $2$-eigenspace is only $1$-dimensional, so $A$ is not
diagonalizable.
\end{solution}



\problem %Dan's notes
Find $\mat{2 6; 0 -1}^{100}$.

\begin{solution}
  $A$ is $2 \times 2$ with $2$ different eigenvalues $\lambda = 2$ and
  $\lambda = -1$, so $A$ is diagonalizable.
  \[\amat{A-2I 0} = \amat{0 6 0; 0 -3 0} \xrightarrow[]{RREF}
    \amat{0 1 0; 0 0 0}, \quad 2\textup{-eigenspace}
    = \Span \left\{\vec{1 0} \right\}.\]
  \[\amat{A+I 0} = \amat{3 6 0; 0 0 0} \xrightarrow[]{RREF}
    \amat{1 2 0; 0 0 0}, \quad (-1)\textup{-eigenspace}
    = \Span \left\{\vec{-2 1} \right\}.\]
  Therefore, $A = CDC^{-1}$ where $C = \mat{1 -2; 0 1}$ and
  $D = \mat{2 0; 0 -1}$.  Using the repeated cancellations of $C$ with
  $C^{-1}$ when expanding $A^{100}$, we get 
  \begin{align*}
    A^{100} & = (C D C^{-1})^{100} = C D^{100} C^{-1} \\
            & = C \mat{2^{100} 0; 0 (-1)^{100}} C^{-1} \\
            & = \mat{1 -2; 0 1} \mat{2^{100} 0; 0 1} \mat{1 2; 0 1} \\
            & = \mat{1 -2; 0 1} \mat{2^{100} 2^{101}; 0 1} \\
            & = \mat{2^{100} 2^{101}-2; 0 1}.
  \end{align*}

\end{solution}

\problem % Jankowski Spring 2018 exam 3  
Let $A = \mat{-1 0 -2; 0 2 0; 3 0 4}$.
\begin{subproblems}
  \subproblem Find the eigenvalues of $A$. %1, 2, 2
  \subproblem Find a basis for each eigenspace of $A$.  Mark your
  answers clearly.
   % $\vec{1 0 1}$, $\vec{0 1 0}$, $\vec{3 0 2}$
  \subproblem \enspace Is $A$ diagonalizable?  If your answer is yes,
  find a diagonal matrix $D$ and an invertible matrix $P$ so that
  $A = PDP^{-1}$.  If your answer is no, justify why $A$ is not
  diagonalizable.
\end{subproblems}

\begin{solution}
  \begin{subproblems}
    \subproblem We solve $0=\det(A-\lambda I)$.
    \begin{align*}
      0 & = \det\mat{-1-\lambda, 0 -2; 0 2-\lambda, 0; 3 0 4-\lambda} =
          (2-\lambda)(-1)^{4} \det\mat{-1-\lambda, -2; 3 4-\lambda} 
      \\ & =(2-\lambda)\left((-1-\lambda)(4-\lambda) + 6\right)
           = (2-\lambda)(\lambda^2-3\lambda -4+6) 
      \\ & = (2-\lambda)(\lambda^2-3\lambda+2)
           = (2-\lambda)(\lambda-2)(\lambda-1)
           % = -(2-\lambda)^2(\lambda-5).
    \end{align*}
    So $\lambda = 1$ and $\lambda = 2$ are the eigenvalues.\newline
    \newline \underline{$\lambda = 1$}:
    $\amat{A-I 0} = \amat{-2 0 -2 0; 0 1 0 0; 3 0 3 0}
    \xrightarrow[\textup{then } R_1 = -R_1
    /2]{R_3=R_3+\frac{3}{2}R_1}\amat{1 0 1 0; 0 1 0 0; 0 0 0 0}$ with
    solution $x_1= -x_3$, $x_2=0$, $x_3=x_3$.  The $1$-eigenspace has
    basis $\left\{ \vec{-1 0 1} \right\}$.  \newline \newline
    \underline{$\lambda = 2$}:
      \[\amat{A-2I 0} = \amat{-3 0 -2 0; 0 0 0 0; 3 0 2 0} 
        \xrightarrow[\textup{then } R_1= -R_1/3]{R_3=R_3+R_1} \amat{1
          0 \frac{2}{3} 0; 0 0 0 0; 0 0 0 0}\] with solution
      $x_1= -\frac{2}{3}x_3$, \ $x_2=x_2$, \ $x_3=x_3$.  The
      $2$-eigenspace has basis
      $\left\{ \vec{0 1 0}, \vec{-2/3 0 1}\right\}$.
      \newline % \newline
      \subproblem $A$ is diagonalizable; $A=PDP^{-1}$ where
      $P = \mat{-1 0 -2/3; 0 1 0; 1 0 1}$ and
      $D= \mat{1 0 0; 0 2 0; 0 0 2}$.
    \end{subproblems}
\end{solution}



\problem %Dan's notes
Let $A = \mat{2 4 6; 0 2 2; 0 0 4}$.  Is $A$ diagonalizable?
If so, find a basis of ${\bf R}^3$ consisting of eigenvectors of $A$.

\begin{solution}
\begin{subproblems}
  \subproblem The matrix is upper-triangular, so the eigenvalues are
  the diagonal entries: $\lambda_1 = -1$, $\lambda_2 = 4$, and
  $\lambda_3 = 5$.
  %
  \subproblem We find a basis for each eigenspace.
  \[\amat{A-\lambda_1I 0} =
    \amat{A+I 0} =
    \amat{0 5 10 0; 0 5 10 0; 0 0 6 0}
    \xrightarrow[]{RREF}
    \amat{0 1 0 0; 0 0 1 0; 0 0 0 0}.\]
  This means $x_2 = x_3 = 0$, so a basis for $(-1)$-eigenspace
  is $\left\{\vec{1 0 0}\right\}$.
  \[\amat{A-\lambda_2I 0} =
    \amat{A-4I 0} =
    \amat{-5 5 10 0; 0 0 10 0; 0 0 1 0}
    \xrightarrow[]{RREF}
    \amat{1 -1 0 0; 0 0 1 0; 0 0 0 0}.\]
  So $x_1 = x_2$ (with $x_2$ free) and $x_3 = 0$, so
  a basis for the $4$-eigenspace is $\left\{\vec{1 1 0}\right\}$.
  %
  \[\amat{A-\lambda_3I 0} =
    \amat{A-5I 0} = \amat{-6 5 10 0; 0 -1 10 0; 0 0 0 0}
    \xrightarrow[]{RREF} \amat{1 0 -10 0; 0 1 -10 0; 0 0 0 0}.\] Thus
  $x_1 = 10x_3$ and $x_2 = 10x_3$ ($x_3$ free). Basis for
  $5$-eigenspace is $\left\{ \vec{10 10 1} \right\}$.
  % 
  \subproblem Yes.  We form $C$ using the three eigenvectors we found
  in part (b), and $D$ is the diagonal matrix with corresponding
  eigenvalues.
  \[C = \mat{1 1 10; 0 1 10; 0 0 1} \qquad
    D = \mat{-1 0 0; 0 4 0; 0 0 5}.\]  
\end{subproblems}
\end{solution}



\problem % Made for ILA
Write a $4 \times 4$ matrix $A$ with only one real eigenvalue $\lambda = 2$, so
that the $2$-eigenspace for $A$ is a three-dimensional subspace of ${\bf
  R}^4$.  %

\begin{solution}
  We need $\lambda=2$ to be the only eigenvalue, and we need $\amat{A-2I, 0}$ to
  have three free variables and thus one pivot.  For example,
    \[A = \mat{2 1 0 0; 0 2 0 0; 0 0 2 0; 0 0 0 2}.\] %
\end{solution}

\problem % Made for ILA
Write a $3 \times 3$ matrix $B$ with the property that $e_1$, $e_2$, and $e_3$
are eigenvectors of $B$.  Must $B$ be diagonalizable?  Justify your answer.

\begin{solution}
Endless possibilities.  For example,
\[B = \mat{1 0 0; 0 1 0; 0 0 1}, \qquad B = \mat{-1 0 0; 0 5 0; 0 0 -3}, \qquad
  B = \mat{2 0 0; 0 3 0; 0 0 3}.\] Note that $B$ must be diagonalizable because
it is a $3 \times 3$ matrix with $3$ linearly independent eigenvectors, namely
$e_1$, $e_2$, $e_3$.  In truth, $B$ satisfies the conditions of this problem if
and only if it is a diagonal $3 \times 3$ matrix.
\end{solution}

  
\problem % Jankowski Spring 2018 Exam 3 (also Sample Exam for Spring
         % 2019)
Find the $2 \times 2$ matrix $A$ whose eigenspaces are drawn below.
Fully simplify your answer.
% \newline
(to be clear: the dashed line is the $(-2)$-eigenspace).
\begin{center}
  \includegraphics[scale=0.5]{e3_pic1}
\end{center}


\begin{solution}
  From the picture, we see $\lambda_1 = 1$ has eigenvector
  $v_1 = \vec{1 2}$.  Also, $\lambda= -2$ has eigenvector
  $v_2 = \vec{-1 1}$.  Forming $P=\mat{v_1 v_2}$ and
  $D=\mat{1 0; 0 -2}$ we get $A=PDP^{-1}$, so
   \begin{align*} A & = \mat{1 -1; 2 1} \mat{1 0; 0 -2} \dfrac{1}{3}
     \mat{1 1; -2 1} \\ & = \frac{1}{3} \mat{1 -1; 2 1} \mat{1 1; 4
       -2} \\ & = \frac{1}{3} \mat{-3 3; 6 0} = \mat{-1 1; 2 0}.
   \end{align*}
\end{solution}

\problem % Created for ILA
Let $A = \mat{1 3; 18 2} \mat{-2 0; 0 3} \mat{1 -5; 18 2}^{-1}$.
\newline Find all values of $k$ (if there are any) so that
\[A \vec{9 k}= \vec{27 3k}.\]
\begin{solution}
  We do not need to do any matrix multiplication, since $A$ has been
  diagonalized for us:
  \begin{itemize}
  \item $\lambda = -2$ is an eigenvalue with corresponding eigenspace
    Span$\left\{\vec{1 18}\right\}$
  \item $\lambda = 3$ is an eigenvalue with corresponding eigenspace
    Span$\left\{\vec{3 2}\right\}$
  \end{itemize}
  The problem is subtly asking us to find $k$ so that $\vec{9 k}$ is
  in the $3$-eigenspace of $A$.  This means $\vec{9 k}$ is a scalar
  multiple of $\vec{3 2}$, so $k = 6$.
\end{solution}

\problem % Made for ILA
\begin{subproblems}
  \subproblem Suppose $A$ is a diagonalizable $n \times n$ matrix, and
  suppose $B$ is an $n \times n$ matrix with the same eigenvalues of
  $A$ and the same corresponding eigenspaces for each eigenvalue.  Is
  it necessarily true that $A=B$? %
  \subproblem Answer part (a) in the case that $A$ is an $n \times n$
  matrix which is \textit{not} diagonalizable.
\end{subproblems}

\begin{solution}
  \begin{subproblems}
    \subproblem Yes.  Since $A$ is diagonalizable, we can write
    $A = PDP^{-1}$, where the columns of $P$ are a basis of
    ${\bf R}^n$ consisting of eigenvectors of $A$ and where $D$ is the
    diagonal matrix with the corresponding eigenvalues.  Since $B$ has
    the same eigenvalues of $A$ and the same corresponding
    eigenspaces, we also have $B = PDP^{-1}$, so $A = B = PDP^{-1}$.
     %
    \subproblem No.  For example, consider $A= \mat{0 1; 0 0}$ and
    $B = \mat{0 2; 0 0}$.  Each matrix has only one eigenvalue
    $\lambda = 0$ with algebraic multiplicity $2$, and the
    $0$-eigenspace for each matrix is the span of $\vec{1 0}$, yet
    $A \neq B$.
  \end{subproblems}
\end{solution}


\problem Consider the following matrices $A$ and $B$.
\[A = \mat{3 1 4; 0 1 2; -1 0 1} \mat{\sqrt{2} 0 0; 0 70^{1/7} 0; 0 0
    \pi} \mat{3 1 4; 0 1 2; -1 0 1}^{-1},\]
\[B = \mat{-3 3 2; 0 3 1; 1 0 \frac{1}{2}} \mat{\sqrt{2} 0 0; 0
    70^{1/7} 0; 0 0 \pi} \mat{-3 3 2; 0 3 1; 1 0 \frac{1}{2}}^{-1}.\]
Without doing any computations, determine whether $A = B$ or
$A \neq B$.

\begin{solution}
  We show that $A=B$.  By the Diagonalization Theorem, $A = PDP^{-1}$,
  where $D$ is any diagonal matrix of eigenvalues of $A$ (with their
  geometric multiplicities) and $P$ is \textit{any} matrix whose
  column vectors form bases for the corresponding eigenspaces.

  The $\sqrt{2}$-eigenspace of $A$ is
  $\textup{Span}\left\{\vec{3 0 -1}\right\}$, which is the same as
  $\textup{Span}\left\{\vec{-3 0 1}\right\}$.

  The $70^{1/7}$-eigenspace of $A$ is $\textup{Span}\vec{1 1 0}$ which
  is the same as $\textup{Span}\left\{\vec{3 3 0} \right\}$.

  The $\pi$-eigenspace of $A$ is
  $\textup{Span}\left\{\vec{4 2 1}\right\}$, which is the same as
  $\textup{Span}\left\{\vec{2 1 \frac{1}{2}}\right\}$.

  Therefore, by the Diagonalization Theorem, $A$ and $B$ are two ways
  of writing the same matrix.
\end{solution}

\problem More practice with diagonalization.  \newline \newline %
Consider the matrix
$A = \begin{pmatrix}-1& 5& 10 \\ 0& 4& 10 \\0 & 0 & 5 \end{pmatrix}$.
\begin{subproblems}
  \subproblem Find the eigenvalues of $A$.

  %\ifanswers \else \vspace{1.5cm} \fi

  \subproblem For each eigenvalue of $A$, compute a basis for the
  corresponding eigenspace.

  %\ifanswers \else \vspace{9cm} \fi

  \subproblem Decide whether $A$ is diagonalizable. If it is
  diagonalizable, find an invertible $3 \times 3$ matrix $P$ and a
  diagonal matrix $D$ such that $A = PDP^{-1}$. If $A$ is not
  diagonalizable, explain why.
\end{subproblems}
\begin{solution} For this problem, we give the answers and leave the
  computations to the reader.
  \begin{subproblems}
    \subproblem The matrix is upper-triangular, so its eigenvalues are
    its diagonal entries: $\lambda = -1$, $\lambda = 4$, and $\lambda = 5$.
    %
    \subproblem 
    The $(-1)$-eigenspace is $\Span\left\{\vec{1 0 0}\right\}$. \newline  %
    The $4$-eigenspace is $\Span\left\{\vec{1 1 0}\right\}$. \newline %
    The $5$-eigenspace is $\Span\left\{\vec{10 10 1}\right\}$.
    %
    \subproblem $A$ is diagonalizable: $A = PDP^{-1}$ where the diagonal entries
    of $D$ are the eigenvalues and the columns of $P$ are eigenvectors written
    in the corresponding order.  Many answers are possible, for example:
    \[P = \mat{1 1 10; 0 1 10; 0 0 1}, \qquad D = \mat{-1 0 0; 0 4 0; 0 0 5}.\]
  \end{subproblems}
\end{solution}

\problem %Dan's notes
More practice with diagonalization.  \newline \newline %
Let $A = \mat{2 0 0; 1 2 1; -1 0 1}$.  Is $A$ diagonalizable?  If so, find a
basis of ${\bf R}^3$ consisting of eigenvectors of $A$.

\begin{solution}
  For this problem, we give the answers and leave the computations to the
  reader. \newline \newline %
  Solving $\det(A-\lambda I) = 0$ gives $\lambda = 1$ and $\lambda = 2$.
  \newline \newline %
  The $1$-eigenspace is $\Span\left\{\vec{0 -1 1}\right\}$. \newline %
  The $2$-eigenspace is %
  $\Span\left\{\vec{0 1 0}, \vec{-1 0 1}\right\}$. \newline %
  The set $\left\{\vec{0 -1 1}, \vec{0 1 0}, \vec{-1 0 1} \right\}$ is a basis
  for ${\bf R}^{3}$ consisting of eigenvectors for $A$, so $A$ is
  diagonalizable.  In fact, $A = PDP^{-1}$ where the diagonal entries of $D$ are
  the eigenvalues and the columns of $P$ are eigenvectors written in the
  corresponding order.  Many answers are possible, for example:
    \[P = \mat{0 0 -1; -1 1 0; 1 0 1}, \qquad D = \mat{1 0 0; 0 2 0; 0 0 2}.\]
\end{solution}


 
\end{problems}\setcounter{problem}{0}

\newpage

\begin{center} \underline{5.5, Complex eigenvalues} \end{center}

{\bf Important note}: In this section, as in the rest of Math 1553, the entries
of all matrices are real numbers unless specified otherwise.

\begin{problems}
\problem % Dan's notes
True or false: If $v$ is an eigenvector for $A$ and $v$ has complex entries,
then $i \cdot v$ is also an eigenvector for $A$.

\begin{solution}
  True: any scalar multiple of a complex eigenvector is also a complex
  eigenvector.  By linearity of matrix multiplication, we see that if
  $Av = \lambda v$, then
  \[A(iv) = i Av= i(\lambda v) = \lambda(iv),\]
  so if $v$ is an eigenvector then so is $iv$.
\end{solution}

\problem % Dan's notes
True or false: If $\vec{i 1}$ is an eigenvector for $A$, then $\vec{i -1}$ is
also an eigenvector for $A$. %tricky but true

\begin{solution}
  True. Recall that the complex conjugate of any complex eigenvector is also a
  complex eigenvector, and any (nonzero) scalar multiple of a complex
  eigenvector is also a complex eigenvector.
  \newline \newline %
  Since $v = \vec{i 1}$ is an eigenvector, so is
  $\overline{v} = \vec{-i 1}$, and so is $-\overline{v} = \vec{i -1}$.
\end{solution}

\problem %Dan's notes
If $A$ is a $3 \times 3$ matrix with real entries, what are the possibilities
for the number of different non-real eigenvalues of $A$?

\begin{solution}
  The maximum is $2$.  Every $3 \times 3$ real matrix has at least one real
  eigenvalue, since its characteristic polynomial is a degree $3$ polynomial
  with real coefficients and therefore must have at least one real root
  (straightforward application of end behavior of odd-degree polynomials and the
  Intermediate Value Theorem).  \newline \newline %
  It is possible for $A$ to have exactly 1 real eigenvalue and $2$ non-real
  eigenvalues.  For example, $A = \mat{0 -1 0; 1 0 0; 0 0 1}$ has eigenvalues
  $1$, $i$, and $-i$.
\end{solution}

\problem %Dan's notes
If $A$ is a $4 \times 4$ matrix with real entries, what are the possibilities
for the number of different non-real eigenvalues of $A$?

\begin{solution}
  The max is $4$.  It is not easy to come up with such a matrix $A$, but it is
  possible.  For example,
  \[A = \mat{0 -1 0 0; 1 0 0 0; 0 0 0 -2; 0 0 2 0}\] has the eigenvalues $i$,
  $-i$, $2i$, and $-2i$.
\end{solution}

\problem % Take problem from exam 2 from Spring 2018, which is also a sample exam from Spring 19
 Consider the matrix
  \[ A = \mat{3 -7; 1 -1} \]
 Find all eigenvalues of $A$.  Simplify your answer.
 For the eigenvalue with negative imaginary part, find an eigenvector.

 \begin{solution}
We compute the characteristic equation:
    \[\begin{split}
        0= \det(A-\lambda I)=(3-\lambda)(-1-\lambda)+7=
        \lambda^2-2\lambda-3+7 = \lambda^2-2\lambda+4.
    \end{split} \]
    By the quadratic formula,
    \[ \lambda = \frac{2\pm\sqrt{4-16}}2 = \dfrac{2 \pm 2i\sqrt{3}}{2}= 1 \pm
      \sqrt{3}i.\]

  Let $\lambda = 1-\sqrt{3}i$.  Then
  \[\amat{A-\lambda I 0} = \amat{2+\sqrt{3}i  -7 0; \star, \star, 0}
    \xrightarrow[R_1 = R_1/(2+\sqrt{3}i)]{\textup{destroy
        $R_2$}} \amat{1 \frac{-7}{2+\sqrt{3}i} 0; 0 0 0} =\amat{1 -2+\sqrt{3}i
      0; 0 0 0}
  \]
  So $x_1 = (2-\sqrt{3}i)x_2$ and $x_2$ is free. An eigenvector is
  $v= \vec{2-\sqrt{3}i 1}$.
  
  
  
  An alternative method for finding an eigenvector, using a trick you may have
  seen in class, is to take the first row $(a \ b)$ of $A-\lambda I_2$ to get an
  eigenvector $\vec{-b a}$:
   %  \[ A - \lambda I_2 = \mat{2+\sqrt{3}i -7; \star, \star}
   %   \quad\implies\quad
   %   v = \vec{7 2+\sqrt{3}i} \]
  \[ A - \lambda I_2 = \mat{2+\sqrt{3}i -7; \star, \star}\] Thus
  $\vec{7 2+\sqrt{3}i}$ is an eigenvector for $\lambda$. This answer is
  equivalent to our answer $v$ above since it is a nonzero scalar multiple of
  $v$, as $\vec{7 2+\sqrt{3}i} = (2+\sqrt{3}i)v$.
 \end{solution}

 \problem %Dan's notes
 Find the complex eigenvalues and a vector in each eigenspace for the following
 matrix.
\[A = \mat{3 -2; 4 -1}.\]

\begin{solution}
The characteristic equation is
\[0 = \det(A-\lambda I) = \lambda^2 - \Tr(A) I + \det(A)
  = \lambda^2 - 2 \lambda + 5,\]
so the eigenvalues are
\[\lambda = \dfrac{2 \pm \sqrt{(-2)^2-5(4)}}{2}
  = \dfrac{2 \pm \sqrt{-16}}{2} = \dfrac{2 \pm 4i}{2} = 1 \pm 2i.\] %
We use the $2 \times 2$ eigenvector trick: if $\lambda$ is an eigenvalue of $A$
and $(a \ b)$ is the first row of $A-\lambda I_2$, then $\vec{-b a}$ is an
eigenvector of $A$ unless $a$ and $b$ are both $0$.  %
\newline %
For $\lambda_1 = 1 + 2i$:
\[\amat{A-(1+2i)I 0} = \amat{2-2i -2 0; (*) (*) 0} = \amat{a b 0; (*) (*) 0}.\]
Therefore, $v_1 = \vec{-b a} = \vec{2 2-2i}$ is an eigenvector for
$\lambda_1 = 1+2i$. \newline \newline %
We do not need to do any work to find an eigenvector for the other eigenvalue:
$v_2 = \overline{v_1} = \vec{2 2+2i}$ is an eigenvector for
$\lambda_2 = \overline{\lambda_1} = 1-2i$.  \newline \newline %
There are many possible answers for $v_1$ and $v_2$, and it is generally
difficult to tell whether vectors are complex scalar multiples of each other.
If you obtained different answers than this solution did for (say) $v_1$, you
can check whether your answer is correct by computing $Av_1$ and seeing if it is
equal to $(1+2i)v_1$.

\end{solution}

\problem %Dan's notes 
Find the complex eigenvalues and eigenvectors for the following matrix.
\[A = \mat{-1 -4; 1 -1}.\]
% Solution needed

\problem %Dan's notes 
Find all eigenvalues and eigenvectors for the following matrix.
\[A = \mat{1 0 -2; 1 3 1; 2 0 1}.\]
% Solution needed

\problem Write a real $2 \times 2$ matrix with eigenvalues $5 \pm 2i$.
% Solution needed
\end{problems}\setcounter{problem}{0}

\newpage

\begin{center} \underline{5.6, Stochastic Matrices} \end{center}

\begin{problems}
\problem % Dan's notes
Is there a stochastic matrix where the $1$-eigenspace has dimension greater than $1$?
% Solution needed

\problem % Dan's notes
Find the steady-state vector for the matrix
\[A = \mat{1/2 1/3; 1/2 2/3}.\]
What vector does $A^n \vec{5 7}$ approach as $n \to \infty$?
% Solution needed


\problem % Dan's notes
Find the steady-state vector for the matrix
\[A = \mat{1/3 1/5 1/4; 1/3 2/5 1/2; 1/3 2/5 1/4}.\]
What vector does $A^n \vec{5 7 38}$ approach as $n \to \infty$?
% Solution needed

\problem % Jankowski Spring 2019 midterm 3
  Consider an internet with four pages: 1, 2, 3, and 4.
\newline
Page 1 links to pages 2 and 3. \newline
Page 2 links to pages 3 and 4. \newline
Page 3 links to pages 1 and 4. \newline
Page 4 links to page 1.
\newline \newline
(i) Write the importance matrix $A$ for this internet.

\noindent (ii) $A$ has a unique steady-state vector $w$.  Find $w$.  You may use the fact
that \[A \vec{8 4 6 5} = \vec{8 4 6 5}.\]

\begin{solution}
  (i) The importance matrix is
    \[A = \mat{0 0 1/2 1; 1/2 0 0 0; 1/2 1/2 0 0; 0 1/2 1/2 0} .\]
    
  (ii) $w=\dfrac{1}{8+4+6+5} \vec{8 4 6 5} = \vec{8/23 4/23 6/23 5/23}.$

\end{solution}

\problem
Courage Soda and Dexter Soda compete for a market of $210$ customers
who drink soda each day.  \newline
Today, Courage has $80$ customers and Dexter has $130$ customers.
Each day:

	$70\%$ of Courage Soda's customers keep drinking 
	Courage Soda, while $30\%$ switch to Dexter Soda.  
	
	$40\%$ of Dexter Soda's customers keep drinking Dexter Soda, while
	$60\%$ switch to Courage Soda.

	\begin{subproblems}
	%\subproblem Write a stochastic matrix $A$ that represents the daily transition of 
	%customers between Courage Soda and Dexter Soda.  		
      \subproblem Write a stochastic matrix $A$ and a vector $x$ so
      that $Ax$ will give the number of customers for Courage Soda and
      Dexter Soda (in that order) tomorrow. \newline
      \ifanswers
        \newline
	
	$A= \mat{0.7 0.6; 0.3 0.4}$ and $x=\vec{80 130}$.
	\newline
	
	\else
	\fi
	
	\subproblem Find the steady-state vector for $A$.
	\ifanswers
	\newline
	
	$\amat{A-I 0}= \amat{-0.3 0.6 0; 0.3 -0.6 0}
    \xrightarrow[R_1=R_1/(-0.3)]{R_2=R_2+R_1} \amat{1 -2 0; 0 0 0}$
    \newline so $x_1 = 2x_2$ and $x_2$ is free.  A $1$-eigenvector is
    $\vec{2 1}$, so the steady state vector is
    $w= \dfrac{1}{2+1} \vec{2 1} = \vec{2/3 1/3}$.  \newline
	
	\else
	\fi
	
	\subproblem Use your answer from (b) to determine the following:
    in the long run, roughly how many daily customers will Courage
    Soda have?
    \ifanswers

      As $n$ gets large, $A^n \vec{80 130}$
      approaches $210 \vec{2/3 1/3} = \vec{140 70}$.  Courage will
      have roughly $140$ customers.  \else \vspace{3cm} \fi
	
	\end{subproblems}

\end{problems}\setcounter{problem}{0}


\newpage